\documentclass[11pt,a4paper]{article}

% ── Packages ──────────────────────────────────────────────────────────
\usepackage[utf8]{inputenc}
\usepackage[T1]{fontenc}
\usepackage{lmodern}
\usepackage{amsmath,amssymb}
\usepackage{booktabs}
\usepackage{enumitem}
\usepackage{hyperref}
\usepackage[margin=1in]{geometry}
\usepackage{xcolor}

\hypersetup{
  colorlinks=true,
  linkcolor=blue!60!black,
  urlcolor=blue!60!black,
}

% ── Title ─────────────────────────────────────────────────────────────
\title{\textbf{OKCA --- Design Document}\\[0.3em]
       \large OK Contrast Algorithm}
\author{}
\date{Status: Production}

\begin{document}
\maketitle
\tableofcontents
\newpage

%% =====================================================================
\section{What OKCA Is Trying to Solve}

WCAG~2.x contrast uses the ratio formula:
\[
  \text{ratio} = \frac{Y_1 + 0.05}{Y_2 + 0.05}
\]

This has two well-documented failure modes:

\begin{enumerate}
  \item \textbf{False passes for saturated chromatic text.}
    Hot pink on near-black scores 6.6:1 under WCAG---a comfortable AA
    pass.  Perceptual research flags it as harder to read than achromatic
    pairs at equivalent luminance.  WCAG's luminance formula does not
    encode the reduced effectiveness of saturated chromatic light on very
    dark backgrounds.

  \item \textbf{False passes for green text/backgrounds near the AA
    boundary.}
    OKLCH lightness cubed, used as a perceptually uniform luminance
    proxy, underestimates WCAG luminance for green hues because
    IEC~61966-2-1 weights sRGB green strongly (0.7152).  A pair that
    WCAG rates 4.4 can appear to score 4.7 when the luminance proxy is
    miscalibrated.
\end{enumerate}

OKCA corrects both, while maintaining \textbf{zero false passes}
against WCAG---meaning OKCA never tells a designer a pair is safe
when WCAG says it fails.

%% =====================================================================
\section{Core Design Principles}

\subsection{\texorpdfstring{Safety: FP\,=\,0 Is Non-Negotiable}{Safety: FP = 0 Is Non-Negotiable}}

A contrast algorithm used for accessibility decisions must never approve
a pair that WCAG rejects.  This is not a calibration preference---it is
the definition of a safe accessibility tool.  OKCA guarantees this
mathematically (see \S\ref{sec:fp0}).

\subsection{WCAG-Compatible Scale and Thresholds}

OKCA outputs ratios in [1,~21], uses the same AA threshold (4.5) and
AAA threshold (7.0) as WCAG~2.x.  Designers and developers see familiar
numbers.  No re-education required.

\subsection{OKLCH L as the Luminance Input}

WCAG~2.x computes luminance from linearised sRGB (IEC~61966-2-1).
OKLCH~L is a perceptually uniform lightness derived from the Oklab
colour space.  For neutral greys, the relationship is straightforward:
\[
  L^3 \approx Y_{\text{WCAG}}
\]

The cube-root perceptual transform inverts to linear luminance.  Using L
as the primary input makes the algorithm coherent with what designers see
in modern colour pickers.

\subsection{Clean-Room IP}

OKCA is a clean-room implementation derived from mathematical
specifications and probe-validated calibration, with no third-party
contrast algorithm source code.

\subsection{Conservative Rather Than Liberal}

When the algorithm must choose between overcounting contrast (false pass)
and undercounting (false failure), it chooses undercounting.  False
failures result in designers choosing safer, more distinct colour pairs.
False passes result in inaccessible text shipped to production.

%% =====================================================================
\section{\texorpdfstring{Why $L^3$ Instead of WCAG Y}{Why L³ Instead of WCAG Y}}

WCAG luminance is defined as:
\[
  Y = 0.2126\,R_{\text{lin}} + 0.7152\,G_{\text{lin}} + 0.0722\,B_{\text{lin}}
\]
per IEC~61966-2-1 (linearised sRGB).  This formula was designed for
display calibration, not for predicting legibility.  Three shortcomings
are relevant to OKCA:

\begin{enumerate}
  \item \textbf{Green is over-weighted.}  The 0.7152 green coefficient
    reflects photopic luminance sensitivity, not perceptual contrast
    effectiveness.  Pure green at Y~=~0.715 scores 14.5:1 against
    black---visually plausible---but formula-level deviations accumulate
    near the AA boundary for dark greens.

  \item \textbf{Not designer-readable.}  A designer cannot look at a
    colour picker, read the OKLCH~L values, and predict the WCAG ratio
    without a separate calculation.

  \item \textbf{Chromatic text on dark backgrounds.}  WCAG~Y does not
    distinguish between saturated and desaturated light text.  A
    chromatic lighter element at the same Y as an achromatic one produces
    the same ratio---even though the saturated one has demonstrably lower
    effective contrast against a dark background.
\end{enumerate}

OKLCH $L^3$ addresses points~2 and~3 directly.  The green
over-weighting (point~1) is corrected by the green correction mechanism
(\S\ref{sec:green}).

%% =====================================================================
\section{\texorpdfstring{The FP\,=\,0 Guarantee}{The FP = 0 Guarantee}}
\label{sec:fp0}

The guarantee rests on two properties of the algorithm, both visible in
the ratio formula:
\[
  \text{ratio} = \frac{Y_{\text{lighter}} + 0.05}{Y_{\text{darker}} + 0.05}
\]

\paragraph{Step~3 (lighter-element compression) can only reduce the
numerator.}

The chroma compression raises L to a power greater than~1 for saturated
colours.  Since $L \le 1$, a higher exponent produces a smaller value:
\[
  L^{\,1.75} \;\le\; L^{\,1.0} \quad \text{for } L \in [0,\,1]
\]
A smaller \texttt{lighterY} in the numerator means a smaller ratio.  For
achromatic colours the exponent is~1 and the value is unchanged.

\paragraph{Step~4 (darker-element green correction) can only increase
the denominator.}

The correction adds a non-negative term to \texttt{darkerL} before
cubing, making \texttt{darkerY} larger.  A larger denominator in the
ratio formula means a lower ratio.  The gate $a < -0.05$ ensures it
fires only where $L^3$ genuinely undershoots WCAG~Y (true greens),
preventing spurious ratio reductions elsewhere.

\medskip
Together, both corrections push the ratio downward:
\[
  \text{ratio}_{\text{OKCA}} \;\le\; \text{ratio}_{\text{WCAG}}
  \quad\text{for any input}
\]
A pair WCAG fails (ratio~<~4.5) will also fail OKCA.
\textbf{FP\,=\,0 by construction.}%

%% =====================================================================
\section{The Two Correction Mechanisms}

\subsection{Chroma Compression on the Lighter Element (Step~3)}
\label{sec:chroma}

\paragraph{Problem.}
Saturated chromatic light text on a dark background---e.g.\ hot pink
(\texttt{\#ff69b4}) on near-black (\texttt{\#1a1a1a})---has the same
WCAG luminance ratio as achromatic white on a dark grey, but is
demonstrably harder to read.  The reduced effectiveness comes from
chromatic contrast partially substituting for luminance contrast in the
visual system.

\paragraph{Mechanism.}
The lighter element's luminance proxy is penalised by a chroma-weighted
power exponent.  First, a saturation weight is computed from the Oklab
chroma:
\[
  \text{satW} = \min\!\left(1,\;\left(\frac{C}{0.15}\right)^{\!2}\right)
\]
This is a quadratic ramp: zero for achromatic colours, reaching 1.0 when
chroma hits the threshold.  The exponent is then:
\[
  \text{exp} = 1 + 0.75 \times \text{satW}
\]
ranging from 1.0 (achromatic) to 1.75 (fully saturated).  The adjusted
luminance proxy becomes:
\[
  Y_{\text{lighter}} = \left(L_{\text{lighter}}^{\;\text{exp}}\right)^3
\]

For a neutral white lighter element: C~=~0, satW~=~0, exp~=~1, so
$Y_{\text{lighter}} = L^3 = 1.0$---unchanged.

For hot pink (C~$\approx$~0.197~>~0.15): satW~=~1, exp~=~1.75---the
lighter element is penalised and the ratio drops below WCAG.

The threshold of 0.15 is calibrated so the penalty is negligible for
lightly tinted neutrals (e.g.\ off-white at C~$\approx$~0.01) and fully
active for vivid designer palette colours (C~$\ge$~0.15).

\subsection{Green Correction on the Darker Element (Step~4)}
\label{sec:green}

\paragraph{Problem.}
For green hues, $L^3$ underestimates WCAG~Y.  Specifically,
\texttt{\#228b22} (forest green) has L~=~0.558, $L^3 = 0.174$, while
WCAG~Y~=~0.189.  Because the denominator in the ratio formula is too
small (darker element appears darker than WCAG says), the ratio is
inflated---creating a false-pass risk.

Root cause: IEC~61966-2-1 assigns 71.5\% weight to the sRGB green
channel.  OKLCH~L's perceptual uniformity distributes weight more
evenly.  The gap is hue-dependent and largest for vivid greens.

\paragraph{Mechanism.}
For the darker element, if its Oklab \texttt{a} coordinate is below the
gate threshold (true greens; not blues, which have a slightly negative
\texttt{a} but at much smaller magnitude), the effective lightness is
boosted:
\[
  L_{\text{eff}} = L_{\text{darker}} +
    \max\!\big(0,\; K_{\text{DARK}} \times (-a - A_{\text{THRESH}})\big)
\]
\[
  Y_{\text{darker}} = L_{\text{eff}}^{\;3}
\]
With $K_{\text{DARK}} = 0.155$ and $A_{\text{THRESH}} = 0.05$.

For forest green (\texttt{\#228b22}): a~=~$-0.135$, correction
=~$0.155 \times (0.135 - 0.05) = 0.013$, so $L_{\text{eff}}$~=~0.558~+~0.013~=~0.571,
and $Y_{\text{darker}}$~=~0.186.  The inflated denominator brings the
ratio down from 4.5 to 4.4---matching WCAG's 4.4.

\paragraph{The gate threshold matters.}
Without it, the correction would fire for slightly blue-shifted colours
like Tailwind blue-700 (\texttt{\#1d4ed8}, a~$\approx$~$-0.047$), where
$L^3 \approx Y_{\text{WCAG}}$ and no correction is needed.  Setting the
gate at 0.05 restricts the correction to colours where the gap is real
and meaningful.

\paragraph{Rounding interaction.}
$K_{\text{DARK}} = 0.155$ gives a raw ratio of 4.448 for
white/\texttt{\#228b22}, which rounds to 4.4 via \texttt{toFixed(1)}.
Probe calibration must account for the production rounding mode---a raw
ratio of 4.471 would round to ``4.5'', producing a false pass at display
precision.

%% =====================================================================
\section{Achromatic Behaviour}

For neutral grey pairs, both corrections reduce to identity:

\begin{itemize}
  \item \textbf{Step~3:} Oklab chroma C~=~0, so satW~=~0, exp~=~1, and
    $Y_{\text{lighter}} = L^3$.
  \item \textbf{Step~4:} Oklab a~=~0, so the gate is not reached and
    $Y_{\text{darker}} = L^3$.
\end{itemize}

Since $L^3 \approx Y_{\text{WCAG}}$ for greys (up to floating-point
precision of the OKLCH transform), the algorithm produces \textbf{exact
WCAG~2.x ratios for achromatic pairs}---white/black~=~21,
white/\texttt{\#767676}~=~4.5.  This is the most important compatibility
property: the AA boundary grey anchor is preserved exactly.

%% =====================================================================
\section{Symmetry}

OKCA is \textbf{symmetric}: okca(A,~B)~=~okca(B,~A).  This follows
from:

\begin{enumerate}
  \item The lighter/darker split uses OKLCH~L magnitude, not polarity
    (fg/bg roles are irrelevant).
  \item Both correction functions depend only on the colour, not on
    whether it is text or background.
  \item The ratio formula uses positional roles determined by~L:
    \[
      \frac{Y_{\text{lighter}} + 0.05}{Y_{\text{darker}} + 0.05}
    \]
\end{enumerate}

Some contrast algorithms use asymmetric polarity models---they model the
difference between reading light text on dark vs.\ dark text on light
backgrounds.  OKCA does not encode this asymmetry.  This is a deliberate
design choice: it produces consistent numbers regardless of which colour
a designer designates as ``text,'' and avoids the calibration complexity
of polarity-dependent response curves.

The cost: OKCA does not capture the finding that warm chromatic light
text on dark backgrounds is harder to read than dark text on light
backgrounds at the same luminance ratio.  It handles this partially via
the chroma compression penalty (step~3), but not via a full polarity
model.

%% =====================================================================
\section{Probe Validation Summary}

Three independent batteries:

\begin{table}[h]
\centering
\begin{tabular}{@{}lrccp{6cm}@{}}
\toprule
\textbf{Battery} & \textbf{Pairs} & \textbf{FP} & \textbf{FF} & \textbf{Notes} \\
\midrule
Light-on-dark   & 53    & \textbf{0} & 1  & Hot pink---intentional WCAG FP \\
Dark-on-light   & 54    & \textbf{0} & 0  & Clean sweep \\
Design systems  & 2{,}480 & \textbf{0} & 28 & All warm-hue conservatism \\
\midrule
\textbf{Total}  & \textbf{2{,}587} & \textbf{0} & \textbf{29} & \\
\bottomrule
\end{tabular}
\end{table}

The 28 design-system false failures are all warm saturated families (red,
fuchsia, pink, rose, orange, plum, indigo) in Tailwind, Material, and
Radix~UI palettes.  They represent principled conservatism: $L^3 >
Y_{\text{WCAG}}$ for warm hues, so OKCA underestimates those pairs
relative to WCAG.  No correction is applied because any warm-side
correction on the lighter element risks false passes for pink/fuchsia
text near the AA boundary.

%% =====================================================================
\section{Key Constants}

\begin{table}[h]
\centering
\begin{tabular}{@{}lrl@{}}
\toprule
\textbf{Constant} & \textbf{Value} & \textbf{Role} \\
\midrule
\texttt{C\_THRESH}  & 0.15  & Oklab chroma at which lighter-element penalty is fully active \\
\texttt{CHROMA\_K}  & 0.75  & Maximum additional power exponent at full saturation \\
\texttt{K\_DARK}    & 0.155 & Green correction coefficient on darker element \\
\texttt{A\_THRESH}  & 0.05  & Oklab \texttt{a} gate: green correction fires only when a~<~$-0.05$ \\
\bottomrule
\end{tabular}
\end{table}

All four constants have one degree of freedom each.  They were calibrated
by the following anchors:

\begin{description}[style=nextline,leftmargin=2em]
  \item[C\_THRESH = 0.15]
    Typical Oklab chroma for designer palette saturated colours.
    Lightly tinted neutrals (C < 0.05) receive less than 10\% of
    the full penalty.

  \item[CHROMA\_K = 0.75]
    Yields exp~=~1.75 at full saturation, which reduces hot pink's
    effective ratio from 6.6 to 4.0 (below the 4.5 AA threshold).  This
    is the target: WCAG's two most-cited false passes (hot pink and dark
    orange on near-black) become OKCA false failures.

  \item[A\_THRESH = 0.05]
    Separates true greens (a~$\approx$~$-0.13$ for forest green) from
    blue-shifted colours like Tailwind blue-700
    (a~$\approx$~$-0.047$).  The gate must be above $-0.047$ and below
    $-0.085$ (the weakest green that needs correction).

  \item[K\_DARK = 0.155]
    Minimum value such that white/\texttt{\#228b22} gives a raw ratio
    $\le$~4.45 (rounds to 4.4, not 4.5).  Derived analytically from the
    target $Y_{\text{darker}} \ge 0.18596$.
\end{description}

%% =====================================================================
\section{What OKCA Does Not Do}

Understanding the scope prevents incorrect use and misguided extension
attempts.

\paragraph{Does not model polarity.}
Reading direction (light-on-dark vs.\ dark-on-light) is not encoded.
Polarity models add significant calibration complexity and require
asymmetric response functions.  OKCA prioritises simplicity and symmetry.

\paragraph{Does not model font size or weight.}
WCAG~AA (4.5:1) applies uniformly regardless of text size.  OKCA outputs
a single ratio; size-dependent thresholds are the caller's
responsibility.

\paragraph{Does not replace perceptual judgement.}
An OKCA score of 4.5 on a warm fuchsia text/white background means the
pair clears the numerical threshold.  A designer may still find it
unpleasant.  OKCA is a safety floor, not a design recommendation.

\paragraph{Does not correct warm-hue false passes in WCAG.}
WCAG~2.x has known false passes for warm chromatic text on very dark
backgrounds.  OKCA catches the two most-cited examples (hot pink, dark
orange on near-black).  It does not systematically correct all warm-hue
WCAG false passes, because doing so would require a polarity model and
risks introducing new false passes.

%% =====================================================================
\section{Extension Guidelines}

When modifying the algorithm, these properties must be preserved:

\begin{enumerate}
  \item \textbf{FP\,=\,0.}  Run all three probe batteries after any
    constant change.  A false pass in the production service requires
    that the probe test be run with rounding applied
    (\texttt{toFixed(1)} on the raw ratio), not just with raw floats.

  \item \textbf{Achromatic exactness.}  White/black~=~21,
    white/\texttt{\#767676}~=~4.5.  These are the WCAG~2.x anchors
    that users will cross-check.  Any deviation breaks trust.

  \item \textbf{No third-party contrast algorithm source code.}
    Mathematical derivations from public specifications are permitted;
    copying control flow or constant blocks from third-party packages is
    not.

  \item \textbf{Update probe battery expectations.}  If any probe result
    changes intentionally, document the rationale before committing.
\end{enumerate}

\end{document}
